\chapter[FORMATS DE RAPPORTS DE POSITION ET RADIOGONIOMÈTRIE (DF)]{Formats de rapports de position et
radiogoniomètrie (DF)}\label{formats-de-rapports-de-position-et-radiogoniomètrie (DF)}

\newfontfamily{\cheightspacesymbolfont}{FreeSans}
\newcommand{\cheightspace}[1]{%
  \begingroup
  \strut{\cheightspacesymbolfont #1}%
  \endgroup
}
\newenvironment{cheightformatbox}[1]{%
  \par\medskip\begingroup
  \setlength{\arrayrulewidth}{0.6pt}%
  \begin{center}%
  \begin{tabular}{|p{0.92\textwidth}|}%
  \hline
  \cellcolor{gray!20}\hspace{0.6em}\sffamily\bfseries\itshape #1\\%
  \hline
  \begin{minipage}{0.98\linewidth}%
  \vspace{0.35em}%
}{%
  \vspace{0.35em}%
  \end{minipage}\\%
  \hline
  \end{tabular}%
  \end{center}%
  \endgroup\par\medskip
}
\newcommand{\chEightHead}[1]{\shortstack{\sffamily\bfseries\itshape #1}}
\newcommand{\chEightPositionHead}[1]{%
  \parbox[c][50pt][c]{\linewidth}{\centering\sffamily\bfseries\itshape #1\par}%
}
\newcommand{\chEightDataExtensionHead}{%
  \makebox[\linewidth][c]{%
  \begin{tabular}{@{}>{\centering\arraybackslash\sffamily\bfseries\itshape}p{\dimexpr\linewidth+2\tabcolsep\relax}@{}}
  Course/Speed \\
  \hline
  Power/Height/Gain/Dir \\
  \hline
  Radio Range \\
  \hline
  DF Signal Strength
  \end{tabular}}%
}
\newcommand{\chEightProto}[1]{\aprsorigtexttt{#1}}
\newcommand{\chEightFitProto}[1]{\resizebox{\linewidth}{!}{\aprsorigtexttt{#1}}}
\newcommand{\chEightBytes}[1]{%
  \makebox[0pt][r]{Bytes:\hspace{1.2em}}%
  \makebox[\linewidth][c]{#1}%
}
\newcommand{\chEightExampleText}[1]{{\rmfamily\small #1}}

\phantomsection\addcontentsline{toc}{section}{Position Reports}
\begin{aprssideblock}{Position Reports}
Les Lat/Long Position Reports sont contenus dans le champ Information
d'une trame AX.25 APRS.

Les diagrammes suivants montrent les formats permis de ces rapports,
accompagnés de quelques exemples. Les zones grises indiquent les champs
optionnels, et les caractères ombrés (jaunes) sont des caractères ASCII
littéraux. Dans tous les cas, il y a un maximum de 43 caractères après
le Symbol Code.

Pourquoi ? D'où vient cette limite de 43 caractères ? Le chapitre 3
indiquait que le champ Information pouvait contenir jusqu'à 256
caractères. En pratique, on voit souvent des commentaires bien plus
longs. Le format Mic-E utilise « n » pour la longueur du commentaire
plutôt qu'une limite précise.
\end{aprssideblock}

\begingroup
\setlength{\arrayrulewidth}{0.6pt}
\begin{center}
\sffamily\fontsize{8.4}{9.4}\selectfont
\setlength{\tabcolsep}{2pt}
\renewcommand{\arraystretch}{1.5}
\begin{tabular}{|>{\centering\arraybackslash}m{0.08\linewidth}|>{\centering\arraybackslash}m{0.08\linewidth}|>{\centering\arraybackslash}m{0.10\linewidth}|>{\centering\arraybackslash}m{0.10\linewidth}|>{\centering\arraybackslash}m{0.12\linewidth}|>{\centering\arraybackslash}m{0.44\linewidth}|}
\hline
\multicolumn{6}{|>{\columncolor{gray!20}}l|}{\hspace{0.6em}\sffamily\bfseries\itshape Format de rapport de position Lat/Long --- sans timestamp}\\
\hline
\chEightPositionHead{\lit{!} or \lit{=}} & \chEightPositionHead{Lat} &
\chEightPositionHead{Sym\\Table\\ID} & \chEightPositionHead{Long} &
\chEightPositionHead{Symbol\\Code} &
\cellcolor{gray!20}\chEightPositionHead{Comment\\(max 43 chars)} \\
\hline
\chEightBytes{1} & 8 & 1 & 9 & 1 & 0--43 \\
\hline
\multicolumn{6}{|p{0.96\linewidth}|}{%
  \begin{minipage}{0.96\linewidth}
  \vspace{0.25em}
  {\rmfamily\underline{Exemples}}

  \vspace{0.25em}
  \begin{tabular}{@{}>{\raggedright\arraybackslash}p{0.50\linewidth}>{\raggedright\arraybackslash}p{0.46\linewidth}@{}}
  \mbox{\chEightProto{!4903.50N/07201.75W-Test 001234}} &
  \chEightExampleText{sans timestamp ni APRS messaging, commentaire.} \\
  \mbox{\chEightProto{!4903.50N/07201.75W-Test /A=001234}} &
  \chEightExampleText{sans timestamp ni APRS messaging, altitude=1234 ft.} \\
  \mbox{\chEightProto{!49}\cheightspace{⎵⎵}\chEightProto{.}\cheightspace{⎵⎵}\chEightProto{N/072}\cheightspace{⎵⎵}\chEightProto{.}\cheightspace{⎵⎵}\chEightProto{W-}} &
  \chEightExampleText{sans timestamp, sans APRS messaging, position au degré le plus proche.} \\
  \chEightProto{!4903.50N/07201.75Wn} &
  \chEightExampleText{sans timestamp, sans APRS messaging.} \\
  \end{tabular}
  \vspace{0.25em}
  \end{minipage}%
} \\
\hline
\end{tabular}%
\end{center}
\endgroup

\begingroup
\setlength{\arrayrulewidth}{0.6pt}
\begin{center}
\sffamily\fontsize{8.4}{9.4}\selectfont
\setlength{\tabcolsep}{2pt}
\renewcommand{\arraystretch}{1.5}
\begin{tabular}{|>{\centering\arraybackslash}m{0.07\linewidth}|>{\centering\arraybackslash}m{0.11\linewidth}|>{\centering\arraybackslash}m{0.07\linewidth}|>{\centering\arraybackslash}m{0.09\linewidth}|>{\centering\arraybackslash}m{0.09\linewidth}|>{\centering\arraybackslash}m{0.11\linewidth}|>{\centering\arraybackslash}m{0.36\linewidth}|}
\hline
\multicolumn{7}{|>{\columncolor{gray!20}}l|}{\hspace{0.6em}\sffamily\bfseries\itshape Format de rapport de position Lat/Long --- avec timestamp}\\
\hline
\chEightPositionHead{\lit{/} or \lit{@}} &
\chEightPositionHead{Time\\DHM /\\HMS} & \chEightPositionHead{Lat} &
\chEightPositionHead{Sym\\Table\\ID} & \chEightPositionHead{Long} &
\chEightPositionHead{Symbol\\Code} &
\cellcolor{gray!20}\chEightPositionHead{Comment\\(max 43 chars)} \\
\hline
\chEightBytes{1} & 7 & 8 & 1 & 9 & 1 & 0--43 \\
\hline
\multicolumn{7}{|p{0.96\linewidth}|}{%
  \begin{minipage}{0.96\linewidth}
  \vspace{0.25em}
  {\rmfamily\underline{Exemples}}

  \vspace{0.25em}
  \begin{tabular}{@{}>{\raggedright\arraybackslash}p{0.49\linewidth}>{\raggedright\arraybackslash}p{0.47\linewidth}@{}}
  \mbox{\chEightProto{/092345z4903.50N/07201.75W>Test1234}} &
  \chEightExampleText{avec timestamp, sans APRS messaging, heure zulu, avec commentaire.} \\
  \mbox{\chEightProto{@092345/4903.50N/07201.75W>Test1234}} &
  \chEightExampleText{avec timestamp, avec APRS messaging, heure locale, avec commentaire.} \\
  \end{tabular}
  \vspace{0.25em}
  \end{minipage}%
} \\
\hline
\end{tabular}%
\end{center}
\endgroup

\begingroup
\setlength{\arrayrulewidth}{0.6pt}
\begin{center}
\sffamily\fontsize{8.0}{9.0}\selectfont
\setlength{\tabcolsep}{3pt}
\renewcommand{\arraystretch}{1.45}
\begin{tabular}{|>{\centering\arraybackslash}m{\dimexpr0.05\textwidth+2pt\relax}|>{\centering\arraybackslash}m{0.05\textwidth}|>{\centering\arraybackslash}m{0.06\textwidth}|>{\centering\arraybackslash}m{0.06\textwidth}|>{\centering\arraybackslash}m{0.075\textwidth}|>{\centering\arraybackslash}m{0.18\textwidth}|>{\centering\arraybackslash}m{\dimexpr0.445\textwidth-14\tabcolsep-8\arrayrulewidth-2pt\relax}|}
\hline
\multicolumn{7}{|>{\columncolor{gray!20}}l|}{\hspace{0.6em}\bfseries\itshape Format de rapport de position Lat/Long --- avec Data Extension (sans timestamp)}\\
\hline
\chEightHead{\lit{!} or \lit{=}} & \chEightHead{Lat} & \chEightHead{Sym\\Table\\ID} & \chEightHead{Long} & \chEightHead{Symbol\\Code} &
\chEightDataExtensionHead &
\cellcolor{gray!20}\chEightHead{Comment\\(max 36 chars)} \\
\hline
\chEightBytes{1} & 8 & 1 & 9 & 1 & 7 & 0--36 \\
\hline
\multicolumn{7}{|p{\dimexpr0.92\textwidth-2\tabcolsep-2\arrayrulewidth\relax}|}{%
  \begin{minipage}{\linewidth}
  \vspace{0.25em}
  {\rmfamily\underline{Exemple}}

  \vspace{0.25em}
  \begin{tabular}{@{}>{\raggedright\arraybackslash}p{0.61\linewidth}>{\raggedright\arraybackslash}p{0.35\linewidth}@{}}
  \chEightProto{=4903.50N/07201.75W\#PHG5132} &
  \chEightExampleText{sans timestamp, avec APRS messaging, avec PHG.} \\
  \chEightFitProto{=4903.50N/07201.75W\lit{\_}225/000g000t050r000p001…h00b10138dU2k} &
  \chEightExampleText{\lit{weather report.}} \\
  \end{tabular}
  \vspace{0.25em}
  \end{minipage}%
} \\
\hline
\end{tabular}
\end{center}
\endgroup

\begingroup
\setlength{\arrayrulewidth}{0.6pt}
\begin{center}
\sffamily\fontsize{7.8}{8.8}\selectfont
\setlength{\tabcolsep}{3pt}
\renewcommand{\arraystretch}{1.45}
\begin{tabular}{|>{\centering\arraybackslash}m{\dimexpr0.045\textwidth+4pt\relax}|>{\centering\arraybackslash}m{0.06\textwidth}|>{\centering\arraybackslash}m{0.045\textwidth}|>{\centering\arraybackslash}m{0.06\textwidth}|>{\centering\arraybackslash}m{0.06\textwidth}|>{\centering\arraybackslash}m{0.07\textwidth}|>{\centering\arraybackslash}m{0.18\textwidth}|>{\centering\arraybackslash}m{\dimexpr0.40\textwidth-16\tabcolsep-9\arrayrulewidth-4pt\relax}|}
\hline
\multicolumn{8}{|>{\columncolor{gray!20}}l|}{\hspace{0.6em}\bfseries\itshape Format de rapport de position Lat/Long --- avec Data Extension et timestamp}\\
\hline
\chEightHead{\lit{/} or \lit{@}} & \chEightHead{Time\\DHM /\\HMS} & \chEightHead{Lat} & \chEightHead{Sym\\Table\\ID} & \chEightHead{Long} & \chEightHead{Symbol\\Code} &
\chEightDataExtensionHead &
\cellcolor{gray!20}\chEightHead{Comment\\(max 36 chars)} \\
\hline
\chEightBytes{1} & 7 & 8 & 1 & 9 & 1 & 7 & 0--36 \\
\hline
\multicolumn{8}{|p{\dimexpr0.92\textwidth-2\tabcolsep-2\arrayrulewidth\relax}|}{%
  \begin{minipage}{\linewidth}
  \vspace{0.25em}
  {\rmfamily\underline{Exemples}}

  \vspace{0.25em}
  \begin{tabular}{@{}>{\raggedright\arraybackslash}p{0.61\linewidth}>{\raggedright\arraybackslash}p{0.35\linewidth}@{}}
  \chEightProto{@092345/4903.50N/07201.75W>088/036} &
  \chEightExampleText{avec timestamp, avec APRS messaging, heure locale, course/speed.} \\
  \chEightProto{@234517h4903.50N/07201.75W>PHG5132} &
  \chEightExampleText{avec timestamp, APRS messaging, heure hours/mins/secs, PHG.} \\
  \chEightProto{@092345z4903.50N/07201.75W>RNG0050} &
  \chEightExampleText{avec timestamp, APRS messaging, heure zulu, radio range.} \\
  \chEightProto{/234517h4903.50N/07201.75W>DFS2360} &
  \chEightExampleText{avec timestamp, heure hours/mins/secs, DF, sans APRS messaging.} \\
  \chEightFitProto{@092345z4903.50N/07201.75W\lit{\_}090/000g000t066r000p000…dUII} &
  \chEightExampleText{\lit{weather report.}} \\
  \end{tabular}
  \vspace{0.25em}
  \end{minipage}%
} \\
\hline
\end{tabular}
\end{center}
\endgroup

\par\medskip
\begingroup
\setlength{\arrayrulewidth}{0.6pt}
\sffamily\fontsize{8.4}{9.4}\selectfont
\setlength{\tabcolsep}{5pt}
\renewcommand{\arraystretch}{1.5}
\noindent\hspace{0.04\textwidth}%
\begin{tabular}{|>{\centering\arraybackslash}m{0.04\textwidth}|>{\centering\arraybackslash}m{0.14\textwidth}|>{\centering\arraybackslash}m{0.04\textwidth}|>{\centering\arraybackslash}m{\dimexpr0.28\textwidth-8\tabcolsep-5\arrayrulewidth\relax}|}
\hline
\multicolumn{4}{|>{\columncolor{gray!20}}c|}{\bfseries\itshape Balise Maidenhead Locator}\\
\hline
\chEightPositionHead{\lit{[}} & \chEightPositionHead{Grid Locator} &
\chEightPositionHead{\lit{]}} & \cellcolor{gray!20}\chEightPositionHead{Comment} \\
\hline
\chEightBytes{1} & 4 or 6 & 1 & n \\
\hline
\multicolumn{4}{|p{\dimexpr0.50\textwidth-2\tabcolsep-2\arrayrulewidth\relax}|}{%
  \vspace{0.25em}
  {\rmfamily\underline{Exemples}}

  \vspace{0.25em}
  \chEightProto{[IO91SX] 35 miles NNW of London}

  \chEightProto{[IO91]}
  \vspace{0.25em}
} \\
\hline
\end{tabular}
\endgroup
\par\medskip

\begin{aprsbodyblock}
L'envoi de données GPS brutes est déconseillé. Il s'agissait d'un
bricolage pour les tout premiers trackers, aux ressources de calcul
insuffisantes pour permettre la conversion en un rapport de position
correct. Les symboles devaient être placés dans le champ destination en
utilisant des noms tels que \texttt{GPSxxx}.
\end{aprsbodyblock}

\par\medskip
\noindent\begin{minipage}{\textwidth}
\vspace*{2em}
\begingroup
\setlength{\arrayrulewidth}{0.6pt}
\sffamily\fontsize{8.4}{9.4}\selectfont
\setlength{\tabcolsep}{0pt}
\renewcommand{\arraystretch}{1.8}
\noindent\hspace{0.04\textwidth}%
\begin{tabular}{|>{\centering\arraybackslash}p{0.05\textwidth}|>{\centering\arraybackslash}p{0.25\textwidth}|>{\centering\arraybackslash}p{0.62\textwidth}|}
\hhline{|-|-~|}
\multicolumn{2}{|>{\columncolor{gray!20}}p{0.30\textwidth}|}{\hspace{0.6em}\bfseries\itshape Format de rapport de position NMEA brut} &
\multicolumn{1}{p{0.62\textwidth}}{} \\
\cline{1-2}
\multicolumn{2}{|>{\centering\arraybackslash}p{0.30\textwidth}|}{\chEightHead{NMEA Received Sentence}} &
\multicolumn{1}{p{0.62\textwidth}}{} \\
\cline{1-2}
\chEightHead{\lit{\$}} & \chEightHead{…,…,…,…,…,…,…} & \multicolumn{1}{p{0.62\textwidth}}{} \\
\cline{1-2}
\chEightBytes{1} & 25--209 & \multicolumn{1}{p{0.62\textwidth}}{} \\
\hline
\multicolumn{3}{|p{0.92\textwidth}|}{%
  \vspace{0.25em}
  \hspace{0.6em}{\rmfamily\scriptsize\underline{Exemples}}

  \vspace{0.25em}
  {\small\ttfamily
  \hspace{0.6em}\$GPGGA,102705,5157.9762,N,00029.3256,W,1,04,2.0,75.7,M,47.6,M,,*62\par
  \hspace{0.6em}\$GPGLL,2554.459,N,08020.187,W,154027.281,A\par
  \hspace{0.6em}\$GPRMC,063909,A,3349.4302,N,11700.3721,W,43.022,89.3,291099,13.6,E*52\par
  \hspace{0.6em}\$GPVTG,318.7,T,,M,35.1,N,65.0,K*69}
  \vspace{0.25em}
} \\
\hline
\end{tabular}
\endgroup
\end{minipage}
\par\vspace{1em}

\phantomsection\addcontentsline{toc}{section}{DF Reports}
\begin{aprssideblock}{DF Reports}
Les DF Reports sont contenus dans le champ Information d'une trame AX.25
APRS. Les paramètres Bearing and Number/Range/Quality (BRG/NRQ) suivent
le champ Data Extension.

\textbf{Note :} les paramètres BRG/NRQ ne sont significatifs que lorsque
le rapport contient le symbole DF (c'est-à-dire que le Symbol Table ID
est \lit{/} et le Symbol Code est \lit{\textbackslash}).

\textbf{Note :} si la station DF est fixe, la valeur Course est zéro. Si
la station est en mouvement, la valeur Course est non nulle.
\end{aprssideblock}

\begingroup
\setlength{\arrayrulewidth}{0.6pt}
\begin{center}
\sffamily\fontsize{7.7}{8.7}\selectfont
\setlength{\tabcolsep}{3pt}
\renewcommand{\arraystretch}{1.45}
\begin{tabular}{|>{\centering\arraybackslash}m{\dimexpr0.045\textwidth+4pt\relax}|>{\centering\arraybackslash}m{0.05\textwidth}|>{\centering\arraybackslash}m{0.06\textwidth}|>{\centering\arraybackslash}m{0.06\textwidth}|>{\centering\arraybackslash}m{0.07\textwidth}|>{\centering\arraybackslash}m{0.18\textwidth}|>{\centering\arraybackslash}m{0.11\textwidth}|>{\centering\arraybackslash}m{\dimexpr0.345\textwidth-16\tabcolsep-9\arrayrulewidth-4pt\relax}|}
\hline
\multicolumn{8}{|>{\columncolor{gray!20}}l|}{\hspace{0.6em}\bfseries\itshape Format de rapport DF --- sans timestamp}\\
\hline
\chEightHead{\lit{!} or \lit{=}} & \chEightHead{Lat} & \chEightHead{Sym\\Table\\ID\\\lit{/}} & \chEightHead{Long} & \chEightHead{Symbol\\Code\\\lit{\textbackslash}} &
\chEightDataExtensionHead &
\chEightHead{\lit{/}BRG\lit{/}NRQ} &
\cellcolor{gray!20}\chEightHead{Comment\\(max 28 chars)} \\
\hline
\chEightBytes{1} & 8 & 1 & 9 & 1 & 7 & 8 & 0--28 \\
\hline
\multicolumn{8}{|p{\dimexpr0.92\textwidth-2\tabcolsep-2\arrayrulewidth\relax}|}{%
  \begin{minipage}{\linewidth}
  \vspace{0.25em}
  {\rmfamily\underline{Exemples}}

  \vspace{0.25em}
  \begin{tabular}{@{}>{\raggedright\arraybackslash}p{0.61\linewidth}>{\raggedright\arraybackslash}p{0.35\linewidth}@{}}
  \chEightProto{=4903.50N/07201.75W\textbackslash088/036/270/729} &
  \chEightExampleText{sans timestamp, course/speed/bearing/NRQ, avec APRS messaging. Station DF en mouvement (CSE est non nul).} \\
  \chEightProto{=4903.50N/07201.75W\textbackslash}\lit{000}\chEightProto{/036/270/729} &
  \chEightExampleText{même rapport, station DF fixe (CSE=\lit{000}).} \\
  \end{tabular}
  \vspace{0.25em}
  \end{minipage}%
} \\
\hline
\end{tabular}
\end{center}
\endgroup

\begingroup
\setlength{\arrayrulewidth}{0.6pt}
\begin{center}
\sffamily\fontsize{7.5}{8.5}\selectfont
\setlength{\tabcolsep}{3pt}
\renewcommand{\arraystretch}{1.45}
\begin{tabular}{|>{\centering\arraybackslash}m{\dimexpr0.04\textwidth+6pt\relax}|>{\centering\arraybackslash}m{0.055\textwidth}|>{\centering\arraybackslash}m{0.04\textwidth}|>{\centering\arraybackslash}m{0.055\textwidth}|>{\centering\arraybackslash}m{0.055\textwidth}|>{\centering\arraybackslash}m{0.06\textwidth}|>{\centering\arraybackslash}m{0.17\textwidth}|>{\centering\arraybackslash}m{0.10\textwidth}|>{\centering\arraybackslash}m{\dimexpr0.345\textwidth-18\tabcolsep-10\arrayrulewidth-6pt\relax}|}
\hline
\multicolumn{9}{|>{\columncolor{gray!20}}l|}{\hspace{0.6em}\bfseries\itshape Format de rapport DF --- avec timestamp}\\
\hline
\chEightHead{\lit{/} or \lit{@}} & \chEightHead{Time\\DHM /\\HMS} & \chEightHead{Lat} & \chEightHead{Sym\\Table\\ID\\\lit{/}} & \chEightHead{Long} & \chEightHead{Symbol\\Code\\\lit{\textbackslash}} &
\chEightDataExtensionHead &
\chEightHead{\lit{/}BRG\lit{/}NRQ} &
\cellcolor{gray!20}\chEightHead{Comment\\(max 28 chars)} \\
\hline
\chEightBytes{1} & 7 & 8 & 1 & 9 & 1 & 7 & 8 & 0--28 \\
\hline
\multicolumn{9}{|p{\dimexpr0.92\textwidth-2\tabcolsep-2\arrayrulewidth\relax}|}{%
  \begin{minipage}{\linewidth}
  \vspace{0.25em}
  {\rmfamily\underline{Exemples}}

  \vspace{0.25em}
  \begin{tabular}{@{}>{\raggedright\arraybackslash}p{0.61\linewidth}>{\raggedright\arraybackslash}p{0.35\linewidth}@{}}
  \chEightProto{@092345z4903.50N/07201.75W\textbackslash088/036/270/729} &
  \chEightExampleText{avec timestamp, course/speed/bearing/NRQ, avec APRS messaging.} \\
  \chEightProto{/092345z4903.50N/07201.75W\textbackslash000/000/270/729} &
  \chEightExampleText{avec timestamp, bearing/NRQ, sans course/speed, sans APRS messaging.} \\
  \end{tabular}
  \vspace{0.25em}
  \end{minipage}%
} \\
\hline
\end{tabular}
\end{center}
\endgroup
